\subsection{公式的形成规则}
\label{formation_rules}

有了 $\mathcal{L}_0$ 的初始符号 \ref{variables_and_connectives}，下面定义 $\mathcal{L}_0$ 的语法规则。符合语法规则的字符串，称为 $\mathcal{L}_0$ 的“公式”，通常记为 $\varphi_0,\varphi_1,\cdots$ 或 $\varphi,\psi$。由公式组成的集合，称为“公式集”，通常记为 $\Gamma_0,\Gamma_1,\cdots$ 或 $\Gamma,\Delta,\Sigma$。

\begin{definition}[命题逻辑公式形成规则]
    \label{formation_rules_of_L0}
    \index{命题逻辑公式形成规则}
    $\mathcal{L}_0$ 公式的形成规则如下：
    \begin{enumerate}
        \item 任意命题变元 $p_0,p_1,\cdots$ 都是 $\mathcal{L}_0$ 的公式
        \item 如果 $\varphi$ 是公式，那么 $\neg \varphi$ 也是公式
        \item 如果 $\varphi$ 和 $\psi$ 是公式，那么 $(\varphi \to \psi)$ 也是公式
    \end{enumerate}
\end{definition}

\begin{definition}[命题逻辑的公式 Formula]
    \label{formula_of_L0}
    \index{命题逻辑的公式 Formula}
    $\varphi$ 是 $\mathcal{L}_0$ 的公式，当且仅当，$\varphi$ 是有限次使用上述形成规则得到的字符串。除此之外没有别的公式。
\end{definition}

\begin{example}
    以下是 $\mathcal{L}_0$ 的一些公式：
    \begin{enumerate}
        \item $p$
        \item $\neg\neg p$
        \item $\neg (p \to q)$
        \item $((p \to q) \to r)$
        \item $\neg (p \to \neg q)$
        \item $(\neg p \to q)$
    \end{enumerate}
\end{example}

\begin{example}
    以下不是 $\mathcal{L}_0$ 的公式：
    \begin{enumerate}
        \item $\neg$
        \item $pq$
        \item $p\to\to q$
        \item $p \to$
        \item $\neg \to p$
        \item $\neg (p \to)$
        \item $\neg (p \to q$
        \item $\neg (p \to q))$
        \item $\neg \neg \to p$
    \end{enumerate}
\end{example}